\documentclass[UTF8]{ctexart}
\usepackage{xeCJK}
\usepackage{geometry}
\geometry{left=2cm,right=2cm,top=2.5cm,bottom=2.5cm}
\setlength{\headheight}{12.7pt}

\usepackage{titlesec}
\titleformat{\section}
  {\normalfont\large\bfseries}
  {\thesection}
  {1em}
  {}
\titlespacing*{\section}
  {0pt}
  {3.5ex plus 1ex minus .2ex}
  {2.3ex plus .2ex}

\usepackage{amsmath}
\usepackage{amssymb}
\usepackage{amsthm}
\usepackage{enumitem}
\setlist[enumerate]{itemsep=0pt,topsep=0pt,parsep=0pt,leftmargin=0pt,itemindent=2.5em}
\setlist[itemize]{itemsep=0pt,topsep=0pt,parsep=0pt,leftmargin=0pt,itemindent=2.5em}
\usepackage{fancyhdr}
\pagestyle{fancy}
\fancyhf{}
\usepackage{graphicx}
\usepackage{hyperref}
\usepackage{indentfirst}
\usepackage{lmodern}

\begin{document}
\setlength{\abovedisplayskip}{1pt}
\setlength{\belowdisplayskip}{1pt}

\section{自然数的乘法}

\hypertarget{sec:定义1.1}{}
\noindent\textbf{定义1.1 (自然数的乘法)}

给定\href{../01 自然数/01 自然数#nameddest=sec:定义1.3}{自然数} $m$,
递归定义映射$\mathrm{mul}_m:\omega\to\omega$如下:
\begin{enumerate}
    \item $\mathrm{mul}_m(0)\overset{\mathrm{def}}{=}0$,
    \item $\mathrm{mul}_m(n^+)\overset{\mathrm{def}}{=}\mathrm{mul}_m(n)+m$.
\end{enumerate}

接下来, 对任意的自然数$m,n$, 定义
\[
    m\cdot n\overset{\mathrm{def}}{=}\mathrm{mul}_m(n),\\[-3pt]
\]
称$\cdot$为\textbf{乘法}, 称$m\cdot n$为$m,n$的\textbf{积}, 也可写作$mn$.

\vspace{10pt}

\hypertarget{sec:定理1.1}{}
\noindent\textbf{定理1.1 (自然数半环)}

自然数集上的乘法有以下性质, 从而是交换幺半环:
\begin{enumerate}
    \item 分配律: $\forall m,n,k\in\omega:\,m(n+k)=mn+mk$,
    \item 结合律: $\forall m,n,k\in\omega:\,(mn)k=m(nk)$,
    \item 单位元: $\forall n\in\omega:\,n1=1n=n$,
    \item 交换律: $\forall m,n\in\omega:\,mn=nm$.
\end{enumerate}

\noindent{\kaishu 证明.}

\begin{enumerate}
    \item 任给自然数$m,n$, 对$k$归纳证明$m(n+k)=mn+mk$.
    
    \hspace{2.17em}
    首先$m(n+0)=mn=mn+0=mn+m0$;
    
    \hspace{2.17em}
    其次任取自然数$k$, 假设$m(n+k)=mn+mk$, 那么
    \[
        m(n+k^+)=m(n+k)^+=m(n+k)+m=mn+mk+m=mn+mk^+.
    \]

    \item 任给自然数$m,n$, 对$k$归纳证明$(mn)k=m(nk)$.
    
    \hspace{2.17em}
    首先$(mn)0=0=m0=m(n0)$; 其次任取自然数$k$, 假设$(mn)k=m(nk)$, 那么
    \[
        (mn)k^+=(mn)k+mn=m(nk)+mn=m(nk+n)=m(nk^+).
    \]

    \item 显然$\forall n\in\omega:\,n1=n$. 再对$n$归纳证明$1n=n$.
    
    \hspace{2.17em}
    首先$1\cdot 0=0$; 其次任取自然数$n$, 假设$1n=n$, 那么
    \[
        1n^+=1n+1=n+1=n^+.
    \]

    \item 第一, 对$m$归纳证明$0m=0$. 证明是显然的.    
    
    \hspace{2.17em}
    第二, 对$m$归纳证明$\forall n\in\omega:\,nm+m=n^+m$.
    
    \hspace{2.17em}
    首先对任意的自然数$n$, 有$n0+0=0=n^+0$;
    
    \hspace{2.17em}
    其次任取自然数$m$,
    假设$\forall n\in\omega:\,nm+m=n^+m$, 那么对任意的自然数$n$有
    \[
        nm^++m^+=nm+n+m+1=nm+m+n+1=n^+m+n^+=n^+m^+.
    \]

    \hspace{2.17em}
    最后任给自然数$m$, 对$n$归纳证明$mn=nm$.
    
    \hspace{2.17em}
    首先$m0=0=0m$; 其次任取自然数$n$, 假设$mn=nm$, 那么
    \begin{equation*}
        mn^+=mn+m=nm+m=n^+m.
    \tag*{\qedsymbol}\end{equation*}
\end{enumerate}





\newpage

\hypertarget{sec:定理1.2}{}
\noindent\textbf{定理1.2}

任给自然数$m,n$, 如果$mn=0$, 那么$m=0$或$n=0$.

\noindent{\kaishu 证明.}

任给非零自然数$m$, 对$n$归纳证明$n\neq 0\to mn\neq 0$.

首先$0=0$, 所以有$0\neq 0\to m0\neq 0$;

其次任取自然数$n$, 假设$n\neq 0\to mn\neq 0$, 那么
$mn^+=mn+m\geqslant m$, 从而$mn^+\neq 0$.

综上, 对任意非零的自然数$m,n$, 都有$mn\neq 0$. \hfill $\qedsymbol$

\vspace{10pt}

最后定义自然数的带余除法.

\vspace{10pt}

\hypertarget{sec:定义1.2}{}
\noindent\textbf{定义1.2 (自然数的带余除法)}

给定自然数$m,n$, 如果$m\neq 0$, 那么存在唯一的一对自然数$p$和$q$使得
\[
    n=mp+q\quad\land\quad q\in m,
\]
分别记作$n/m$和$n\%m$, 称作$n$除以$m$的\textbf{商}和\textbf{余数}.

\noindent{\kaishu 证明.}

\noindent{\kaishu 存在性.}

任给非零自然数$m$, 对$n$归纳证明$\exists p,q\in\omega:\,n=mp+q\land q<m$.

首先显然$0=m0+0$并且$0<m$;

其次任取自然数$n$, 假设存在自然数$p,q$使得$n=mp+q$且$q<m$.

此时$q^+<m^+$, 即$q^+<m$或$q^+=m$.

\begin{enumerate}
    \item 如果$q^+<m$, 那么$n^+=mp+q^+$并且$q^+<m$.
    \item 如果$q^+=m$, 那么$n^+=mp+m=mp^+=mp^++0$并且$0<m$.
\end{enumerate}
综上, 存在自然数$p',q'$使得$n^+=mp'+q'$并且$q'<m$.

\noindent{\kaishu 唯一性.}

假设同时存在自然数$p,q$和$p',q'$使得
\[
    n=mp+q=mp'+q'\quad\land\quad q<m\quad\land\quad q'<m.
\]
假设$p<p'$, 即$p^+\leqslant p'$, 则
\[
    n=mp+q<mp+m=mp^+\leqslant mp'\leqslant mp'+q'=n,
\]
矛盾. 同理$p'<p$也矛盾, 所以$p=p'$. 进一步显然$q=q'$. \hfill $\qedsymbol$

\end{document}
